\chapter{СУДЛАГДСАН БАЙДЛЫН ТОЙМ}

Дадал үүсэх үйл явц болон зан үйлийг дэмжих дижитал системүүдийн талаарх судалгааны анхаарал сүүлийн жилүүдэд нэмэгдэж байна. Ялангуяа гар утсанд суурилсан хэрэглээний програм, мэдэгдлийн систем, ахицыг хянах хэрэгсэл, тоглоомжуулалтын элемент, нөхцөл байдалд суурилсан дэмжлэг зэрэг шийдлүүдийг ашиглан хэрэглэгчийн давтагдах зан үйлийг дэмжих оролдлогууд өргөжсөн байна. Энэ чиглэлийн судалгаанууд нь нэг талаас дадал үүсэх сэтгэлзүйн болон зан үйлийн механизмыг тайлбарлах, нөгөө талаас тэдгээрийг дижитал орчинд хэрхэн хэрэгжүүлж болохыг тодорхойлоход чиглэж байна. Иймээс энэ бүлэгт дадал үүсэх онолын үндэс, зан үйлийг дэмжих дижитал системүүдийн талаарх өмнөх болон сүүлийн үеийн судалгааг авч үзэж, улмаар уг ажлын судалгааны цоорхойг тодорхойлно.

\section{Дадал үүсэх механизмын талаарх судалгаа}

Дадал үүсэх механизмыг тайлбарлах судалгааны суурь нь зан үйлийн сэтгэл судлалын эртний туршилтуудаас эхтэй. Павловын сонгодог нөхцөлдөлтийн туршилтаар анх саармаг байсан дохио давталтын үр дүнд автомат хариу үйлдэл өдөөж болдгийг харуулсан \cite{pavlov1927}. Харин Скиннерийн оперант нөхцөлдөлтийн судалгаагаар эерэг үр дагавартай үйлдэл дахин давтагдах магадлал нэмэгддэг болохыг тогтоосон \cite{skinner1938}. Эдгээр судалгаанууд нь зан үйл нь орчны өдөөгч дохио болон түүний дараах үр дагавартай уялдан тогтворжиж болдгийг харуулсан анхны онолын суурь болсон бөгөөд орчин үеийн дадлын судалгаанд өдөөгч дохио--хариу үйлдлийн холбоо болон зан үйлийг бэхжүүлэх бататгалын механизмын үндсэн ойлголт болж байна.

Дадал нь тодорхой нөхцөлд харьцангуй автоматаар гүйцэтгэгддэг зан үйлийн хэлбэр бөгөөд олон удаагийн давталт, орчны тогтвортой нөхцөл, өдөөгч дохионуудын нөлөөгөөр аажмаар бүрэлддэг \cite{ref1}. Wood ба R\"unger нарын тайлбарласнаар дадал нь орчны дохио болон давтагдах зан үйлийн хооронд үүссэн автомат холбоо бөгөөд хүний зан үйл нь зөвхөн ухамсартай шийдвэр гаргалтаар бус, давтагдсан нөхцөлд автоматаар идэвхждэг хэвшмэл хариу үйлдлээр ч ихээхэн тодорхойлогддог \cite{ref1}.

Lally нар шинэ дадал үүсэх үйл явцыг бодит амьдралын нөхцөлд судлахдаа зан үйл тогтмол давтагдах тусам автоматжилтын түвшин өсдөг боловч энэ өсөлт шугаман бус болохыг тогтоосон \cite{ref6}. Эхний үед автоматжилт харьцангуй хурдан өсөх хандлагатай байдаг бол тодорхой түвшинд хүрсний дараа аажмаар тогтворжин, тэгширсэн төлөвт ордог. Тэдний судалгаанд автоматжилтын өсөлтийг асимптотик муруйгаар тайлбарласан бөгөөд 95 хувийн дадал тогтворжих медиан хугацаа ойролцоогоор 66 өдөр, хувь хүний түвшинд 18--254 хоногийн хэлбэлзэлтэй байжээ \cite{ref6}.

Иймээс дадал үүсэх үйл явц нь нэг удаагийн шийдвэр, богино хугацааны сэдэлжилтээс илүүтэйгээр өдөөгч дохио, хариу үйлдэл, давталт, бататгал, улмаар автоматжилт гэсэн урт хугацааны уялдаанд суурилдаг гэж үзэж болно.

\section{Дадлын хүчний ойлголт ба хэмжилтийн талаарх судалгаа}

Дадлын хүчийг судалгаанд ихэвчлэн тухайн зан үйл хэр зэрэг автоматжсан бэ гэдгээр ойлгодог. Өөрөөр хэлбэл, дадлын хүч нь зөвхөн гүйцэтгэлийн тоо, тасралтгүй гүйцэтгэлийн цуваа, эсвэл өндөр гүйцэтгэлийн хувиар тодорхойлогдохгүй. Судалгааны ажлуудад дадлын гол механизм нь автоматжилт гэж тайлбарладаг бөгөөд давтамж нь автоматжилтын шалтгаан болон үр дагаврын аль аль нь байж болох ч дадлын хүчний шууд хэмжүүр биш гэж үздэг \cite{ref1}.

Дадлын хүчийг хэмжихэд өргөн хэрэглэгддэг хэмжүүрүүдийн нэг нь SRHI (Self-Report Habit Index) бөгөөд энэ нь давталт, автоматжилт, өөрийн дүр төрхтэй холбогдох байдал зэрэг хэд хэдэн талыг хамардаг \cite{ref8}. Харин SRBAI (Self-Report Behavioral Automaticity Index) нь SRHI-ээс гаргаж авсан, автоматжилтын хэмжээс дээр төвлөрсөн богино хувилбар юм \cite{ref9}. Сүүлийн үеийн хэмжилзүйн судалгаанд SRBAI нь дадлын хүчийг тооцоолоход илүү парсимони, ойлгомжтой, өргөн хэрэглэгдэх боломжтой хэмжүүр гэж үзэж байна \cite{ref9}.

Мөн дадлын хүч нь цаг хугацааны явцад шугаман байдлаар өсдөггүй. Автоматжилт нь зан үйл тогтмол давтагдах тусам өсөх боловч тодорхой түвшинд хүрсний дараа аажмаар тогтворжих хандлагатай байдаг \cite{ref6}. Иймээс дадлын хүчийг нэг удаагийн хэмжилтээр бус, хугацааны явц дахь өөрчлөлтийн чиг хандлагаар авч үзэх нь илүү оновчтой юм.

\section{Дадал төлөвшлийн мөчлөг ба зан үйлийн загваруудын талаарх судалгаа}

Дадал үүсэх үйл явцыг тайлбарлахад өргөн хэрэглэгддэг загваруудын нэг нь дадал төлөвшлийн мөчлөг юм. Duhigg уг загварыг өдөөгч дохио -- үйлдлийн хэвшил -- урамшуулал гэсэн гурван үе шаттайгаар тайлбарласан \cite{ref2}. Түүний үзэж байгаагаар орчны тодорхой дохио нь зан үйлийг эхлүүлж, үйлдэл гүйцэтгэгдсэний дараа үүсэх урамшууллын мэдрэмж нь тухайн зан үйлийг бататган, улмаар давталтын үр дүнд автоматжуулдаг \cite{ref2}.

Clear уг хандлагыг өдөөгч дохио -- хүсэл тэмүүлэл -- хариу үйлдэл -- урамшуулал гэсэн илүү нарийвчилсан дөрвөн үе шаттай загвар болгон тайлбарласан \cite{ref3}. Энэхүү тайлбар нь зөвхөн өдөөгч ба үйлдлийн холбоог бус, тухайн үйлдлийг хийхэд нөлөөлөх дотоод хүсэл, хүлээлт, сэдлийн хүчин зүйлийг хамруулдгаараа онцлогтой \cite{ref3}. Иймээс дижитал системийн зохиомжид хөрвүүлэхэд илүү тохиромжтой.

Мөн Fogg-ийн зан үйлийн загварын дагуу тодорхой зан үйл хэрэгжихийн тулд сэдэл, боломж, өдөөгч гэсэн гурван хүчин зүйл тухайн мөчид нэгэн зэрэг бүрдсэн байх ёстой \cite{ref4}. Өөрөөр хэлбэл, хэрэглэгч өндөр сэдэлтэй байсан ч тухайн үйлдэл хийхэд төвөгшил өндөр байвал зан үйл хэрэгжихгүй. Үүний эсрэгээр, үйлдэл хийхэд хялбар байсан ч зөв мөчид өдөөгч дохио байхгүй бол мөн хэрэгжилт тасалдаж болно \cite{ref4}.

\section{Дижитал дадал хянах системүүдийн талаарх судалгаа}

Дадлыг дэмжих зорилготой дижитал системүүдийн дийлэнх нь дадал хянах систем хэлбэрээр хөгжсөн байдаг. Эдгээр нь хэрэглэгчид зорилго тодорхойлох, өдөр тутмын гүйцэтгэлийг бүртгэх, тасралтгүй гүйцэтгэлийн цуваа тооцох, сануулга өгөх, ахицыг хянах, статистик мэдээлэл харуулах зэрэг суурь функцүүдийг ашигладаг. Энэ төрлийн системүүд нь зан үйлийн өөрчлөлтийг дэмжихэд тодорхой хувь нэмэр оруулдаг боловч бодит дадал төлөвшилд үзүүлэх нөлөө нь үргэлж тогтвортой байдаггүй \cite{ref5}.

Renfree нарын чанарын судалгаагаар сануулга болон тасралтгүй гүйцэтгэлийн цуваа зэрэг механизмууд нь хэрэглэгчийн давталтыг дэмждэг боловч хэрэглэгчийг системийн сануулгаас хэт хамааралтай болгох эрсдэлтэйг тэмдэглэсэн \cite{ref5}. Өөрөөр хэлбэл, хэрэглэгч тухайн дадлыг дотоод автомат зан үйл болгон төлөвшүүлэхийн оронд системийн дохиогоор өдөөгдөн давтах хандлагатай болж, улмаар бодит дадал үүсэх нөхцөл буурах магадлалтай \cite{ref5}.

Stawarz, Cox, Blandford нарын судалгаанд шинэ дадлыг төлөвшүүлэхдээ жижиг, тодорхой, бага хүчин чармайлт шаардсан үйлдлээс эхлэх нь илүү үр дүнтэй байж болохыг онцолсон \cite{ref11}. Мөн өмнө нь тогтсон зан үйлтэй шинэ дадлыг холбох буюу дадал давхарлах нь тогтвортой байдлыг нэмэгдүүлэх боломжтойг тэмдэглэсэн \cite{ref11}. Энэ нь системийн зүгээс дадлыг жижиг алхамд хуваах, өмнөх хэвшилтэй холбох, хэрэгжүүлэхэд хялбар болгох дэмжлэг үзүүлэх шаардлагатайг харуулж байна.

\section{Тоглоомжуулалт ба урамшууллын системийн талаарх судалгаа}

Хэрэглэгчийн оролцоо, идэвх, тасралтгүй хэрэглээг нэмэгдүүлэхэд тоглоомжуулалт өргөн ашиглагддаг. Hamari, Koivisto, Sarsa нарын тоймоос үзэхэд тоглоомжуулалтын элементүүд нь олон тохиолдолд хэрэглэгчийн идэвх, сэдэл, оролцоог нэмэгдүүлэх эерэг нөлөөтэй байж болох боловч үр дүн нь нөхцөл байдлаас ихээхэн хамаардаг \cite{ref13}.

Гэсэн хэдий ч тоглоомжуулалт нь богино хугацааны идэвхжилд үр дүнтэй байж болох боловч урт хугацааны автомат дадал үүсэхэд дангаараа хангалтгүй байж болно \cite{ref14}. Иймээс урамшууллын зохиомж нь зөвхөн гаднын шагналд тулгуурлах бус, хэрэглэгчийн бодит ахиц, өөрчлөлтийн мэдрэмж, дотоод бататгалтай холбогдсон байх шаардлагатай.

\section{Итгүүлэн нөлөөлөх технологи ба яг цаг тухайд нь дасан зохицсон оролцооны талаарх судалгаа}

Итгүүлэн нөлөөлөх технологи нь хэрэглэгчийг албадахгүйгээр, харин зөөлөн дэмжлэг, санал сэтгэгдэл, өөрийгөө хянах, урамшуулал зэрэг механизмаар дамжуулан зан үйлийн өөрчлөлтийг дэмжих чиглэл юм \cite{ref15}. Oinas-Kukkonen ба Harjumaa нарын боловсруулсан итгүүлэн нөлөөлөх системийн зохиомжийн хүрээ нь ийм төрлийн системийг төлөвлөх, хөгжүүлэх, үнэлэхэд өргөн хэрэглэгддэг онолын суурийн нэг юм \cite{ref15}.

Сүүлийн жилүүдэд энэ чиглэл яг цаг тухайд нь дасан зохицсон оролцоо гэсэн ойлголттой уялдан хөгжиж байна. Энэ төрлийн оролцоо нь хэрэглэгчид зөв төрлийн дэмжлэгийг зөв цагт, зөв нөхцөлд хүргэхийг зорьдог. Ийм оролцоо нь ялангуяа эрүүл мэнд, биеийн идэвхтэй холбоотой зан үйлийг дэмжихэд бага боловч эерэг үр нөлөө үзүүлж болохыг сүүлийн тоймууд харуулж байна \cite{ref16}.

\section{Судалгааны дүгнэлт ба судалгааны цоорхой}

Энэ бүлэгт авч үзсэн судалгаануудыг нэгтгэн дүгнэвэл, дадал нь нөхцөлд суурилсан, давталтаар бэхждэг, автоматаар идэвхждэг зан үйл болох нь тогтоогдож байна \cite{ref1,ref6}. Мөн дижитал зан үйлийн өөрчлөлтийн оролцоо, дадал хянах систем, тоглоомжуулалт, итгүүлэн нөлөөлөх технологи зэрэг чиглэлийн судалгаанууд нь сануулга, ахицыг хянах, өөрийгөө хянах, урамшуулал, нөхцөл байдалд нийцсэн дэмжлэг зэрэг механизмууд тодорхой үр нөлөөтэй болохыг харуулж байна \cite{ref5,ref13,ref15,ref16}.

Гэсэн хэдий ч хэд хэдэн чухал цоорхой ажиглагдаж байна. Нэгдүгээрт, ихэнх дижитал систем өөрийгөө хянах, зорилго тогтоох, сануулга өгөх, ахицыг хянах зэрэг хэсэгчилсэн техникүүдэд төвлөрдөг бөгөөд дадал төлөвшлийн мөчлөгийн бүх үе шатыг системийн функц, өгөгдлийн загвар, алгоритмын түвшинд уялдаатай хэрэгжүүлсэн шийдэл ховор байна. Хоёрдугаарт, дадлын хүчийг үнэлэхдээ олон систем гүйцэтгэлийн давтамж, тасралтгүй гүйцэтгэлийн цуваа зэрэг шууд бус үзүүлэлтүүдэд хэт түшиглэдэг бөгөөд автоматжилтыг тогтмол хэмждэггүй. Гуравдугаарт, сануулга, тоглоомжуулалт, цуваа зэрэг механизм нь богино хугацааны давталтыг дэмждэг ч хэрэглэгчийг системийн дохио, интерфэйсээс хэт хамааралтай болгох эрсдэлтэй байна \cite{ref5}. Дөрөвдүгээрт, онолын түвшинд боловсруулсан зан үйлийн загваруудыг функциональ шаардлага, өгөгдлийн бүтэц, системийн зохиомж, алгоритмын логик, үнэлгээний үзүүлэлт болгон системтэй хөрвүүлсэн инженерчлэлийн судалгаа харьцангуй цөөн байна.

Иймээс энэхүү судалгааны ажил нь дадал төлөвшлийн мөчлөгийн логикийг програм хангамжийн системийн архитектур, өгөгдлийн загвар, системийн зохиомж, үндсэн алгоритмын түвшинд буулгаж, өдөөгч дохио, хүсэл тэмүүлэл, хариу үйлдэл, урамшуулал гэсэн үе шатуудыг уялдуулан дэмжихийн зэрэгцээ дадлын хүчийг автоматжилтад суурилсан хэмжүүрээр үнэлэх туршилтын загвар боловсруулахыг зорьж байна.

